\frfilename{mt261.tex} 
\versiondate{11.12.12} 
\copyrightdate{2000} 
      
\def\dist{\mathop{\text{dist}}} 
      
\def\chaptername{Perubahan variabel dalam integral} 
\def\sectionname{Teorema Vitali dalam $\BbbR^r$} 
      
\newsection{261} 
      
\def\headlinesectionname{Teorema Vitali dalam $\eightBbb R^r$} 
      
Tujuan utama bagian ini ialah memberikan versi berdimensi $r$ dari teorema 
Vitali dan Teorema Kerapatan Lebesgue, dengan mengikuti gagasan-gagasan yang 
telah disajikan dalam \S\S221 dan 223. Saya mengakhirinya dengan pembuktian
bahwa ukuran luar Lebesgue dapat didefinisikan melalui peliputan oleh bola,
bukan oleh interval (261F).
      
\cmmnt{ 
\leader{261A}{Notasi} Dalam sebagian besar bab ini, kita akan berurusan 
dengan geometri dan ukuran ruang Euklides; penetapan sejumlah notasi akan 
menghemat ruang. 
      
Di seluruh bagian ini dan dua bagian berikutnya, $r\ge 1$ adalah suatu 
bilangan bulat. Saya akan menggunakan huruf Latin untuk anggota $\BbbR^r$ dan 
huruf Yunani untuk koordinatnya, sehingga 
$a=(\alpha_1,\ldots,\alpha_r)$, dan seterusnya; jika Anda melihat huruf Yunani 
berindeks bawah, pertama-tama carilah vektor di dekatnya yang mungkin memiliki 
huruf itu sebagai koordinat. Ukuran yang dibahas hampir selalu merupakan 
ukuran Lebesgue pada $\Bbb R^r$; jadi, kecuali dinyatakan lain, $\mu$ harus 
ditafsirkan sebagai ukuran Lebesgue dan $\mu^*$ sebagai ukuran luar Lebesgue. 
Demikian pula, $\int\ldots dx$ selalu berarti integrasi terhadap ukuran 
Lebesgue (dalam dimensi yang ditentukan oleh konteks). 
      
Untuk $x=(\xi_1,\ldots,\xi_r)\in\BbbR^r$, tuliskan 
$\|x\|=\sqrt{\xi_1^2+\ldots+\xi_r^2}$. Ingat bahwa 
$\|x+y\|\le\|x\|+\|y\|$ (1A2C) dan bahwa $\|\alpha x\|=|\alpha|\|x\|$ 
untuk sebarang vektor $x$, $y$ dan skalar $\alpha$. 
      
Saya akan menggunakan notasi yang sama seperti dalam \S115 untuk `interval', 
sehingga, khususnya, 
      
\Centerline{$\coint{a,b} 
=\{x:\alpha_i\le\xi_i<\beta_i\Forall i\le r\}$,} 
      
\Centerline{$\ooint{a,b} 
=\{x:\alpha_i<\xi_i<\beta_i\Forall i\le r\}$,} 
      
\Centerline{$[a,b] 
=\{x:\alpha_i\le \xi_i\le \beta_i\Forall i\le r\}$} 
      
\noindent setiap kali $a$, $b\in\BbbR^r$. 
      
$\tbf{0}=(0,\ldots,0)$ adalah vektor nol dalam $\BbbR^r$, dan 
$\tbf{1}$ adalah $(1,\ldots,1)$. Jika 
$x\in\BbbR^r$ dan $\delta>0$, maka $B(x,\delta)$ adalah bola tertutup 
berpusat di $x$ dan berjari-jari $\delta$, yakni 
$\{y:y\in\BbbR^r,\,\|y-x\|\le\delta\}$. 
Perhatikan bahwa $B(x,\delta)=x+B(\tbf{0},\delta)$; maka, berdasarkan 
invariansi ukuran Lebesgue terhadap translasi, kita mempunyai 
      
\Centerline{$\mu B(x,\delta)=\mu B(\tbf{0},\delta) 
=\beta_r\delta^r$,} 
      
\noindent dengan 
      
$$\eqalign{\beta_r 
&=\Bover{1}{k!}\pi^k\text{ jika }r=2k\text{ genap},\cr 
&=\Bover{2^{2k+1}k!}{(2k+1)!}\pi^k\text{ jika }r=2k+1\text{ ganjil}\cr}$$ 
      
\noindent (252Q). 
}%end of comment 
      
\leader{261B}{Teorema Vitali dalam $\BbbR^r$} Misalkan $A\subseteq\BbbR^r$ 
sebarang himpunan, dan $\Cal I$ suatu keluarga bola tertutup nontrivial\cmmnt{ 
(yakni, bukan singleton, atau secara ekuivalen, tidak terabaikan)} 
dalam $\BbbR^r$ sedemikian sehingga setiap titik dalam $A$ termuat dalam anggota 
$\Cal I$ yang berdiameter sekecil apa pun. Maka terdapat suatu himpunan 
terhitung yang 
saling lepas $\Cal I_0\subseteq \Cal I$ sedemikian sehingga 
$\mu(A\setminus\bigcup\Cal I_0)=0$. 
      
\proof{{\bf (a)} Sebagai permulaan (sampai akhir (f) di bawah), andaikan 
bahwa $\|x\|<M$ untuk setiap $x\in A$, dan tetapkan 
      
\Centerline{$\Cal 
I'=\{I:I\in\Cal I,\,I\subseteq B(\tbf{0},M)\}$.} 
      
\noindent Jika terdapat suatu himpunan berhingga yang saling lepas 
$\Cal I_0\subseteq\Cal I'$ sedemikian sehingga 
$A\subseteq\bigcup\Cal I_0$ (termasuk kemungkinan bahwa 
$A=\Cal I_0=\emptyset$), kita dapat berhenti. Jadi, selanjutnya andaikan 
tidak ada $\Cal I_0$ demikian. 
      
\medskip 
      
{\bf (b)} Dalam hal ini, jika $\Cal I_0$ adalah sebarang subhimpunan berhingga 
yang saling lepas dari $\Cal I'$, terdapat $J\in\Cal I'$ yang saling lepas 
dari setiap anggota $\Cal I_0$. \Prf\ Ambil 
$x\in A\setminus\bigcup\Cal I_0$. Karena setiap anggota $\Cal I_0$ tertutup, 
terdapat $\delta>0$ sedemikian sehingga $B(x,\delta)$ tidak bertemu dengan 
anggota mana pun dari $\Cal I_0$, dan karena $\|x\|<M$ kita dapat mengandaikan 
bahwa $B(x,\delta)\subseteq B(\tbf{0},M)$. Ambil $J$ sebagai anggota 
$\Cal I$ yang memuat $x$ dan berdiameter paling besar $\delta$; maka 
$J\in\Cal I'$ dan $J\cap\bigcup\Cal I_0=\emptyset$.\ 
\Qed 
      
\medskip 
      
{\bf (c)} Karena itu, kita dapat memilih suatu barisan bilangan real 
$\sequencen{\gamma_n}$ dan suatu barisan yang saling lepas 
$\sequencen{I_n}$ dalam $\Cal I'$ secara induktif sebagai berikut. Setelah 
$\langle I_j\rangle_{j<n}$ diberikan (jika $n=0$, ini adalah barisan kosong, 
tanpa anggota), dengan $I_j\in\Cal I'$ untuk setiap $j<n$, dan 
$I_j\cap I_k=\emptyset$ untuk $j<k<n$, tetapkan $\Cal 
J_n=\{I:I\in\Cal I',\,I\cap I_j=\emptyset$ untuk setiap $j<n\}$. Dari (b) 
kita mengetahui bahwa $\Cal J_n\ne\emptyset$. Tetapkan 
      
\Centerline{$\gamma_n=\sup\{\diam I:I\in\Cal J_n\}$;} 
      
\noindent maka $\gamma_n\le 2M$, karena setiap anggota $\Cal J_n$ termuat 
dalam $B(\tbf{0},M)$. Karena itu, kita dapat menemukan suatu himpunan 
$I_n\in\Cal J_n$ sedemikian sehingga $\diam I_n\ge{1\over 2}\gamma_n$, dan 
ini melanjutkan proses induksi. 
      
\medskip 
      
{\bf (d)} Karena $I_n$ merupakan subhimpunan terukur yang saling lepas dari 
himpunan terbatas $B(\tbf{0},M)$, kita mempunyai 
      
\Centerline{$\sum_{n=0}^{\infty}\mu I_n 
\le\mu B(\tbf{0},M)<\infty$,} 
      
\noindent dan $\lim_{n\to\infty}\mu I_n=0$. Juga, 
$\mu I_n\ge\beta_r(\bover14\gamma_n)^r$ untuk setiap $n$, sehingga 
$\lim_{n\to\infty}\gamma_n=0$. 
      
\medskip

{\bf (e)}

Sekarang definisikan $I'_n$ sebagai bola tertutup dengan pusat yang sama 
seperti $I_n$ tetapi dengan diameter lima kali lipat, sehingga bola ini 
memuat setiap titik yang berjarak paling besar $\gamma_n$ dari $I_n$. Saya 
menyatakan bahwa, untuk setiap $n$, 
$A\subseteq\bigcup_{j<n}{I}_j\cup\bigcup_{j\ge n}I_j'$. 
\Prf\Quer\ Andaikan, jika mungkin, sebaliknya. Ambil sebarang 
$x\in A\setminus(\bigcup_{j<n}{I}_j\cup\bigcup_{j\ge n}I'_j)$. Ambil 
$\delta>0$ sedemikian sehingga 
      
      
\Centerline{$B(x,\delta) 
  \subseteq B(\tbf{0},M)\setminus\bigcup_{j<n}I_j$,} 
      
\noindent dan ambil $J\in\Cal I$ sedemikian sehingga 
$x\in J\subseteq B(x,\delta)$. Maka 
      
\Centerline{$\lim_{m\to\infty}\gamma_m=0<\diam J$} 
      
\noindent (di sinilah kita menggunakan hipotesis bahwa semua bola dalam 
$\Cal I$ nontrivial); misalkan $m$ bilangan bulat terkecil yang lebih besar 
atau sama dengan $n$ sedemikian sehingga $\gamma_m<\diam J$. Dalam hal ini 
$J$ tidak mungkin termasuk dalam $\Cal J_m$, sehingga harus ada suatu $k<m$ 
sedemikian sehingga $J\cap I_k\ne\emptyset$, karena tentu saja 
$J\in\Cal I'$. Berdasarkan pemilihan $\delta$, $k$ tidak mungkin kurang dari 
$n$, sehingga $n\le k<m$, dan $\gamma_k\ge\diam J$. Jadi jarak dari $x$ ke 
titik terdekat $I_k$ paling besar $\diam J\le \gamma_k$. Namun, ini berarti 
bahwa $x\in I'_k$; bertentangan dengan pemilihan $x$.\ \Bang\Qed 
      
\medskip 
      
{\bf (f)} Akibatnya, 
      
\Centerline{$\mu\sp*(A\setminus\bigcup_{j<n}I_j)\le\mu(\bigcup_{j\ge 
n}I'_j)\le\sum_{j=n}^{\infty}\mu I'_j\le 5^r\sum_{j=n}^{\infty}\mu 
I_j$.} 
      
\noindent Karena 
      
\Centerline{$\sum_{j=0}^{\infty}\mu I_j\le \mu B(\tbf{0},M)<\infty$,} 
      
\noindent $\lim_{n\to\infty}\mu\sp*(A\setminus\bigcup_{j<n}I_j)=0$ dan 
      
\Centerline{$\mu(A\setminus\bigcup_{j\in\Bbb N}I_j)= 
\mu\sp*(A\setminus\bigcup_{j\in\Bbb N}I_j)=0$.} 
      
      
Jadi, dalam hal ini kita dapat menetapkan $\Cal I_0=\{I_n:n\in\Bbb N\}$ untuk 
memperoleh suatu keluarga terhitung yang saling lepas dalam $\Cal I$ dengan 
$\mu(A\setminus\bigcup\Cal I_0)=0$. 
      
      
\medskip 
      
{\bf (g)} Ini menyelesaikan pembuktian jika $A$ terbatas. Secara umum, 
tetapkan 
      
\Centerline{$U_n=\{x:x\in \BbbR^r,\,n<\|x\|<n+1\}$, 
\quad $A_n=A\cap U_n$, 
\quad$\Cal J_n=\{I:I\in\Cal I,\,I\subseteq U_n\}$,} 
      
\noindent untuk setiap $n\in\Bbb N$. Maka, untuk setiap $n$, kita melihat 
bahwa setiap titik dalam $A_n$ termasuk dalam anggota $\Cal J_n$ yang berdiameter 
sekecil apa pun, sehingga terdapat suatu $\Cal J'_n\subseteq\Cal J_n$ yang terhitung 
dan saling lepas sedemikian sehingga 
$A_n\setminus\bigcup\Cal J'_n$ terabaikan. Sekarang (karena $U_n$ saling 
lepas) $\Cal I_0=\bigcup_{n\in\Bbb N}\Cal J'_n$ saling lepas, dan tentu saja 
$\Cal I_0$ merupakan subhimpunan terhitung dari $\Cal I$; selain itu, 
      
\Centerline{$A\setminus\bigcup\Cal I_0 
\subseteq(\BbbR^r\setminus\bigcup_{n\in\Bbb N}U_n) 
\cup\bigcup_{n\in\Bbb N}(A_n\setminus\bigcup\Cal J'_n)$} 
      
\noindent terabaikan. (Untuk melihat bahwa $\Bbb 
R^r\setminus\bigcup_{n\in\Bbb N}U_n=\{x:\|x\|\in\Bbb N\}$ terabaikan, 
perhatikan bahwa untuk setiap $n\in\Bbb N$ himpunan 
      
\Centerline{$\{x:\|x\|=n\}\subseteq B(\tbf{0},n)\setminus 
B(\tbf{0},\delta n)$} 
      
\noindent mempunyai ukuran paling besar $\beta_rn^r-\beta_r(\delta n)^r$ 
untuk setiap $\delta\in\coint{0,1}$, sehingga harus terabaikan.) 
}%end of proof of 261B 
      
      
\leader{261C}{}\cmmnt{ Seperti dalam \S223, kita dapat menggunakan teorema 
Vitali berdimensi $r$ untuk membuktikan teorema-teorema tentang pendekatan 
fungsi melalui nilai rata-rata lokalnya. 
      
\medskip 
      
\noindent}{\bf Teorema Kerapatan dalam $\BbbR^r$: bentuk integral} Misalkan 
$D$ suatu subhimpunan $\BbbR^r$, dan $f$ suatu fungsi bernilai real yang 
terintegralkan pada $D$. Maka 
      
\Centerline{$f(x)=\lim_{\delta\downarrow 0}\Bover{1}{\mu B(x,\delta)} 
  \int_{D\cap B(x,\delta)}fd\mu$} 
      
\noindent untuk hampir setiap $x\in D$. 
      
\proof{{\bf (a)} Sebagai permulaan (sampai akhir (b)), andaikan bahwa 
$D=\dom f=\BbbR^r$. 
      
Ambil $n\in\Bbb N$ dan $q$, $q'\in\Bbb Q$ dengan $q<q'$, dan tetapkan 
      
\Centerline{$A=A_{nqq'} 
=\{x:\|x\|\le n,\,f(x)\le q,\, 
  \limsup_{\delta\downarrow 0}\Bover{1}{\mu B(x,\delta)} 
  \int_{B(x,\delta)}fd\mu>q'\}$.} 
      
\noindent\Quer\ Andaikan, jika mungkin, bahwa $\mu^*A>0$. Ambil 
$\epsilon>0$ sedemikian sehingga $\epsilon(1+|q|)<(q'-q)\mu^* A$, dan ambil 
$\eta\in\ocint{0,\epsilon}$ sedemikian sehingga $\int_{E}|f|\le\epsilon$ 
setiap kali $\mu E\le\eta$ (225A). Ambil suatu himpunan terbuka $G\supseteq A$ 
yang berukuran paling besar $\mu^*A+\eta$ (134Fa). Misalkan $\Cal I$ adalah 
himpunan bola tertutup nontrivial $B\subseteq G$ sedemikian sehingga 
$\bover{1}{\mu B}\int_Bfd\mu\ge q'$. Maka setiap titik dalam $A$ termuat dalam 
(bahkan merupakan pusat dari) anggota $\Cal I$ yang berdiameter sekecil apa 
pun. Jadi, 
berdasarkan 261B, terdapat suatu himpunan terhitung yang saling lepas 
$\Cal I_0\subseteq\Cal I$ sedemikian sehingga 
$\mu(A\setminus\bigcup\Cal I_0)=0$; tetapkan $H=\bigcup\Cal I_0$. 
      
Karena $\int_{I}fd\mu\ge q'\mu I$ untuk setiap $I\in\Cal I_0$, kita 
mempunyai 
      
\Centerline{$\int_{H}fd\mu=\sum_{I\in\Cal I_0}\int_Ifd\mu 
\ge q'\sum_{I\in\Cal I_0}\mu I=q'\mu H\ge q'\mu^* A$.} 
      
\noindent Tetapkan 
      
\Centerline{$E=\{x:x\in G,\,f(x)\le q\}$.} 
      
\noindent Maka $E$ terukur, dan $A\subseteq E\subseteq G$; jadi 
      
\Centerline{$\mu^* A\le\mu E\le\mu G\le\mu^* A+\eta 
\le\mu^* A+\epsilon$.} 
      
\noindent Selain itu, 
      
\Centerline{$\mu(H\setminus E)\le\mu G-\mu E\le\eta$,} 
      
\noindent sehingga, berdasarkan pemilihan $\eta$, 
$\int_{H\setminus E}f\le\epsilon$ dan 
      
$$\eqalignno{\int_{H}f 
&\le\epsilon+\int_{H\cap E}f 
\le\epsilon+q\mu(H\cap E)\cr 
&\le\epsilon+q\mu^*A+|q|(\mu(H\cap E)-\mu^* A) 
\le q\mu^* A+\epsilon(1+|q|)\cr 
\displaycause{karena $\mu^*A=\mu^*(A\cap H)\le\mu(H\cap E)\le\mu E$} 
&<q'\mu^* A 
\le\int_{H}f,\cr}$$ 
      
\noindent yang mustahil.\ \Bang 
      
Jadi $A_{nqq'}$ terabaikan. Ini berlaku untuk semua $q<q'$ dan semua $n$, 
sehingga 
      
\Centerline{$A^*=\bigcup_{q,q'\in\Bbb Q,q<q'} 
  \bigcup_{n\in\Bbb N}A_{nqq'}$} 
      
\noindent terabaikan. Namun, 
      
\Centerline{$f(x) 
\ge\limsup_{\delta\downarrow 0}\Bover{1}{\mu B(x,\delta)} 
  \int_{B(x,\delta)}f$} 
      
\noindent untuk setiap $x\in\BbbR^r\setminus A^*$, yakni, untuk hampir 
setiap $x\in \BbbR^r$. 
      
\medskip 
      
{\bf (b)} Dengan cara serupa, atau dengan menerapkan hasil ini pada $-f$, 
      
\Centerline{$f(x) 
\le\liminf_{\delta\downarrow 0}\Bover{1}{\mu B(x,\delta)} 
  \int_{B(x,\delta)}f$} 
      
\noindent untuk hampir setiap $x$, sehingga 
      
\Centerline{$f(x) 
=\lim_{\delta\downarrow 0}\Bover{1}{\mu B(x,\delta)} 
  \int_{B(x,\delta)}f$} 
      
\noindent untuk hampir setiap $x$. 
      
\medskip 
      
{\bf (c)} Untuk kasus yang (secara sepintas) lebih umum sebagaimana dinyatakan 
dalam teorema, misalkan $\tilde f$ suatu fungsi yang terintegralkan terhadap 
$\mu$ dan memperluas $f\restr D$, terdefinisi di seluruh $\BbbR^r$, serta 
sedemikian sehingga $\int_F\tilde f=\int_{D\cap F}f$ untuk setiap himpunan 
terukur $F\subseteq\BbbR^r$ (dengan menerapkan 214Eb pada $f\restr D$). Maka 
      
\Centerline{$f(x)=\tilde f(x) 
=\lim_{\delta\downarrow 0}\Bover{1}{\mu B(x,\delta)} 
  \int_{B(x,\delta)}\tilde f 
=\lim_{\delta\downarrow 0}\Bover{1}{\mu B(x,\delta)} 
  \int_{D\cap B(x,\delta)}f$} 
      
\noindent untuk hampir setiap $x\in D$. 
}%end of proof of 261C 
      
\leader{261D}{Akibat} (a) Jika $D\subseteq\BbbR^r$ sebarang himpunan, maka 
      
\Centerline{$\lim_{\delta\downarrow 0}\Bover{\mu^*(D\cap B(x,\delta))} 
{\mu B(x,\delta)}=1$} 
      
\noindent untuk hampir setiap $x\in D$. 
      
(b) Jika $E\subseteq\BbbR^r$ suatu himpunan terukur, maka 
      
\Centerline{$\lim_{\delta\downarrow 0}\Bover{\mu(E\cap B(x,\delta))} 
{\mu B(x,\delta)}=\chi E(x)$} 
      
\noindent untuk hampir setiap $x\in \BbbR^r$. 
      
      
(c) Jika $D\subseteq\BbbR^r$ dan $f:D\to\Bbb R$ sebarang fungsi, maka untuk 
hampir setiap $x\in D$, 
      
\Centerline{$\lim_{\delta\downarrow 0} 
\Bover{\mu^*(\{y:y\in D,\,|f(y)-f(x)|\le\epsilon\}\cap 
B(x,\delta))} 
{\mu B(x,\delta)}=1$} 
      
\noindent untuk setiap $\epsilon>0$. 
      
(d) Jika $D\subseteq\BbbR^r$ dan $f:D\to\Bbb R$ terukur, maka untuk hampir 
setiap $x\in D$, 
      
\Centerline{$\lim_{\delta\downarrow 0} 
\Bover{\mu^*(\{y:y\in D,\,|f(y)-f(x)|\ge\epsilon\}\cap 
B(x,\delta))} 
{\mu B(x,\delta)}=0$} 
      
\noindent untuk setiap $\epsilon>0$. 
      
\proof{{\bf (a)} Terapkan 261C dengan $f=\chi B(\tbf{0},n)$ untuk melihat 
bahwa, untuk setiap $n\in\Bbb N$, 
      
\Centerline{$\lim_{\delta\downarrow 0}\Bover{\mu^*(D\cap B(x,\delta))} 
{\mu B(x,\delta)}=1$} 
      
\noindent untuk hampir setiap $x\in D$ dengan $\|x\|<n$. 
      
\medskip 
      
{\bf (b)} Terapkan (a) pada $E$ untuk melihat bahwa 
      
\Centerline{$\liminf_{\delta\downarrow 0}\Bover{\mu(E\cap 
B(x,\delta))}{\mu B(x,\delta)}\ge\chi E(x)$} 
      
\noindent untuk hampir setiap $x\in\BbbR^r$, dan pada 
$E'=\BbbR^r\setminus E$ untuk melihat bahwa 
      
\Centerline{$\limsup_{\delta\downarrow 0}\Bover{\mu(E\cap 
B(x,\delta))}{\mu B(x,\delta)} 
=1-\liminf_{\delta\downarrow 0} 
  \Bover{\mu(E'\cap B(x,\delta))}{\mu B(x,\delta)} 
\le 1-\chi E'(x)=\chi E(x)$} 
      
\noindent untuk hampir setiap $x$. 
      
\medskip 
      
{\bf (c)} Untuk $q$, $q'\in\Bbb Q$, tetapkan 
      
\Centerline{$D_{qq'}=\{x:x\in D,\,q\le f(x)\le q'\}$,} 
      
\Centerline{$C_{qq'}=\{x:x\in D_{qq'},\,\lim_{\delta\downarrow 0} 
\Bover{\mu^*(D_{qq'}\cap B(x,\delta))}{\mu B(x,\delta)}=1\}$;} 
      
\noindent sekarang tetapkan 
      
\Centerline{$C 
=D\setminus\bigcup_{q,q'\in\Bbb Q}(D_{qq'}\setminus C_{qq'})$,} 
      
\noindent sehingga $D\setminus C$ terabaikan. Jika $x\in C$ dan 
$\epsilon>0$, maka terdapat $q$, $q'\in\Bbb Q$ sedemikian sehingga 
$f(x)-\epsilon\le q\le f(x)\le q'\le f(x)+\epsilon$, dan sekarang 
$x\in C_{qq'}$; oleh karena itu 
      
\Centerline{$\liminf_{\delta\downarrow 0} 
\Bover{\mu^*\{y:y\in D\cap B(x,\delta),\,|f(y)-f(x)|\le\epsilon\}} 
  {\mu B(x,\delta)} 
\ge\liminf_{\delta\downarrow 0} 
  \Bover{\mu^*(D_{qq'}\cap B(x,\delta))}{\mu B(x,\delta)}=1$,} 
      
\noindent sehingga 
      
\Centerline{$\lim_{\delta\downarrow 0} 
\Bover{\mu^*\{y:y\in D\cap B(x,\delta),\,|f(y)-f(x)|\le\epsilon\}} 
{\mu B(x,\delta)} 
=1$.} 
      
\medskip 
      
{\bf (d)} Definisikan $C$ seperti dalam (c). Berdasarkan (a), kita mengetahui 
bahwa $\mu(D\setminus C')=0$, dengan 
      
\Centerline{$C'=\{x:x\in D,\lim_{\delta\downarrow 0}\Bover{\mu^*(D\cap 
B(x,\delta))} 
{\mu B(x,\delta)}=1\}$. 
} 
      
\noindent Jika $x\in C\cap C'$ dan $\epsilon>0$, berdasarkan (c) kita 
mengetahui bahwa 
      
 \Centerline{$\lim_{\delta\downarrow 0} 
\Bover{\mu^*\{y:y\in D\cap B(x,\delta),\,|f(y)-f(x)|\le\epsilon/2\}} 
{\mu B(x,\delta)} 
=1$.} 
      
\noindent Namun, karena $f$ terukur, kita mempunyai 
      
$$\eqalign{\mu^*\{y:y\in &D\cap B(x,\delta),\, 
|f(y)-f(x)|\ge\epsilon\}\cr 
&+\mu^*\{y:y\in D\cap B(x,\delta),\,|f(y)-f(x)|\le\Bover12\epsilon\} 
\le\mu^*(D\cap B(x,\delta))\cr}$$ 
      
\noindent untuk setiap $\delta>0$. Oleh karena itu, 
      
$$\eqalign{\limsup_{\delta\downarrow 0} 
&\Bover{\mu^*\{y:y\in D\cap B(x,\delta),\,|f(y)-f(x)|\ge\epsilon\}} 
{\mu B(x,\delta)}\cr 
&\le\lim_{\delta\downarrow 0} 
  \Bover{\mu^*(D\cap B(x,\delta))}{\mu B(x,\delta)} 
-\lim_{\delta\downarrow 0} 
  \Bover{\mu^*\{y:y\in D\cap B(x,\delta),\,|f(y)-f(x)|\le\epsilon/2\}} 
  {\mu B(x,\delta)} 
=0,\cr}$$ 
      
\noindent dan 
      
\Centerline{$\lim_{\delta\downarrow 0} 
\Bover{\mu^*\{y:y\in D\cap B(x,\delta),\,|f(y)-f(x)|\ge\epsilon\}} 
{\mu B(x,\delta)}=0$} 
      
\noindent untuk setiap $x\in C\cap C'$, yakni, untuk hampir setiap 
$x\in D$. 
}%end of proof of 261D 
      
\leader{261E}{Teorema} Misalkan $f$ suatu fungsi yang terintegralkan secara 
lokal dan terdefinisi pada suatu subhimpunan koterabaikan dari $\BbbR^r$. Maka 
      
\Centerline{$\lim_{\delta\downarrow 0}\Bover1{\mu B(x,\delta)} 
\int_{B(x,\delta)}|f(y)-f(x)|dy=0$} 
      
\noindent untuk hampir setiap $x\in\BbbR^r$. 
      
\proof{ (Bandingkan 223D.) 
      
\medskip 
      
{\bf (a)} Untuk sementara, tetapkan $n\in\Bbb N$, dan tetapkan 
$G=\{x:\|x\|<n\}$. Untuk setiap $q\in\Bbb Q$, tetapkan 
$g_q(x)=|f(x)-q|$ untuk $x\in G\cap\dom f$; maka $g_q$ terintegralkan pada 
$G$, dan 
      
\Centerline{$\lim_{\delta\downarrow 0}\Bover1{\mu B(x,\delta)} 
\int_{G\cap B(x,\delta)}g_q=g_q(x)$} 
      
\noindent untuk hampir setiap $x\in G$, berdasarkan 261C. Dengan menetapkan 
      
\Centerline{$E_q=\{x:x\in G\cap\dom f,\,\lim_{\delta\downarrow 0} 
\Bover1{\mu B(x,\delta)}\int_{G\cap B(x,\delta)}g_q=g_q(x)\}$,} 
      
\noindent kita mempunyai $G\setminus E_q$ terabaikan untuk setiap $q$, 
sehingga $G\setminus E$ terabaikan, dengan 
$E=\bigcap_{q\in\Bbb Q}E_q$. Sekarang 
      
\Centerline{$\lim_{\delta\downarrow 0}\Bover1{\mu B(x,\delta)} 
\int_{G\cap B(x,\delta)}|f(y)-f(x)|dy=0$} 
      
\noindent untuk setiap $x\in E$. \Prf\ Ambil $x\in E$ dan $\epsilon>0$. 
Maka terdapat $q\in\Bbb Q$ sedemikian sehingga $|f(x)-q|\le\epsilon$, 
sehingga 
      
\Centerline{$|f(y)-f(x)|\le|f(y)-q|+\epsilon=g_q(y)+\epsilon$} 
      
\noindent untuk setiap $y\in G\cap\dom f$, dan 
      
$$\eqalign{\limsup_{\delta\downarrow 0}\Bover1{\mu B(x,\delta)} 
\int_{G\cap B(x,\delta)}|f(y)-f(x)|dy 
&\le\limsup_{\delta\downarrow 0}\Bover1{\mu B(x,\delta)} 
\int_{G\cap B(x,\delta)}g_q(y)+\epsilon\,dy\cr 
&=\epsilon+g_q(x) 
\le 2\epsilon.\cr}$$ 
      
\noindent Karena $\epsilon$ sebarang, 
      
\Centerline{$\lim_{\delta\downarrow 0}\Bover1{\mu B(x,\delta)} 
\int_{G\cap B(x,\delta)}|f(y)-f(x)|dy=0$,} 
      
\noindent sebagaimana diperlukan.\ \Qed 
      
\medskip 
      
{\bf (b)} Karena $G$ terbuka, 
 
\Centerline{$\lim_{\delta\downarrow 0}\Bover1{\mu B(x,\delta)} 
\int_{B(x,\delta)}|f(y)-f(x)|dy 
=\lim_{\delta\downarrow 0}\Bover1{\mu B(x,\delta)} 
\int_{G\cap B(x,\delta)}|f(y)-f(x)|dy 
=0$} 
      
\noindent untuk hampir setiap $x\in G$. Karena $n$ sebarang, 
 
\Centerline{$\lim_{\delta\downarrow 0}\Bover1{\mu B(x,\delta)} 
\int_{B(x,\delta)}|f(y)-f(x)|dy=0$} 
      
\noindent untuk hampir setiap $x\in\BbbR^r$. 
}%end of proof of 261E 
      
\medskip 
\noindent{\bf Catatan} Himpunan 
      
\Centerline{$\{x:x\in\dom f,\, 
  \lim_{\delta\downarrow 0}\Bover1{\mu B(x,\delta)} 
\int_{B(x,\delta)}|f(y)-f(x)|dy=0\}$} 
      
\noindent kadang-kadang disebut {\bf himpunan Lebesgue} dari $f$. 
      
\leader{261F}{}\cmmnt{ Akibat lain yang sangat berguna dari 261B adalah 
sebagai berikut. 
      
\medskip 
      
\noindent}{\bf Proposisi} Misalkan $A\subseteq\BbbR^r$ sebarang himpunan, 
dan $\epsilon>0$. Maka terdapat suatu barisan $\sequencen{B_n}$ bola tertutup 
dalam $\BbbR^r$, semuanya berjari-jari paling besar $\epsilon$, sedemikian 
sehingga $A\subseteq\bigcup_{n\in\Bbb N}B_n$ dan 
$\sum_{n=0}^{\infty}\mu B_n\le\mu^*A+\epsilon$. Selain itu, kita dapat 
mengandaikan bahwa bola-bola dalam barisan yang pusatnya tidak terletak dalam 
$A$ memiliki jumlah ukuran paling besar $\epsilon$. 
      
\proof{{\bf (a)} Langkah pertama adalah pengamatan yang jelas bahwa jika 
$x\in\BbbR^r$, $\delta>0$, maka kubus setengah terbuka 
$I=\coint{x,x+\delta\tbf{1}}$ merupakan subhimpunan bola 
$B(x,\delta\sqrt{r})$, yang berukuran 
$\gamma_r\delta^r=\gamma_r\mu I$, dengan 
$\gamma_r=\beta_rr^{r/2}$. Akibatnya, jika $G\subseteq\BbbR^r$ sebarang 
himpunan terbuka, maka $G$ dapat diliput oleh suatu barisan bola dengan jumlah 
ukuran paling besar $\gamma_r\mu G$. \Prf\ Jika $G$ kosong, kita dapat 
mengambil semua bola sebagai himpunan satu titik. Jika tidak, untuk setiap 
$k\in\Bbb N$, tetapkan 
      
\Centerline{$Q_k 
=\{z:z\in \Bbb Z^r,\,\coint{2^{-k}z,2^{-k}(z+\tbf{1})}\subseteq G\}$,} 
      
\Centerline{$E_k 
=\bigcup_{z\in Q_k}\coint{2^{-k}z,2^{-k}(z+\tbf{1})}$.} 
      
\noindent Maka $\sequence{k}{E_k}$ merupakan barisan himpunan yang tak 
menurun dengan gabungan $G$, dan $E_0$ serta setiap selisih 
$E_{k+1}\setminus E_k$ dapat dinyatakan sebagai gabungan yang saling lepas 
dari kubus-kubus setengah terbuka. Dengan demikian, $G$ juga dapat dinyatakan 
sebagai gabungan yang saling lepas dari suatu barisan $\sequencen{I_n}$ kubus 
setengah terbuka. Setiap $I_n$ diliput oleh suatu bola $B_n$ yang berukuran 
$\gamma_r\mu I_n$; sehingga $G\subseteq\bigcup_{n\in\Bbb N}B_n$ dan 
      
\Centerline{$\sum_{n=0}^{\infty}\mu B_n 
\le\gamma_r\sum_{n=0}^{\infty}\mu I_n=\gamma_r\mu G$. \Qed} 
      
\medskip 
      
{\bf (b)} Langsung diperoleh bahwa jika $\mu A=0$, maka untuk setiap 
$\epsilon>0$ terdapat suatu barisan $\sequencen{B_n}$ bola yang meliput $A$ 
dengan jumlah ukuran paling besar $\epsilon$, karena tentu saja terdapat 
suatu himpunan terbuka yang memuat $A$ dan berukuran paling besar 
$\epsilon/\gamma_r$. 
      
\medskip 
      
{\bf (c)} Sekarang ambil sebarang himpunan $A$, dan $\epsilon>0$. Ambil suatu 
himpunan terbuka $G\supseteq A$ dengan 
$\mu G\le\mu^*A+\bover12\epsilon$. Misalkan $\Cal I$ adalah keluarga bola 
tertutup nontrivial yang termuat dalam $G$, berjari-jari paling besar 
$\epsilon$ dan berpusat dalam $A$. Maka setiap titik dalam $A$ termasuk dalam 
anggota $\Cal I$ yang berdiameter sekecil apa pun, sehingga terdapat suatu 
$\Cal I_0\subseteq\Cal I$ yang terhitung dan saling lepas sedemikian sehingga 
$\mu(A\setminus\bigcup\Cal I_0)=0$. Misalkan $\sequencen{B'_n}$ suatu 
barisan bola yang meliput $A\setminus\bigcup\Cal I_0$ dengan 
$\sum_{n=0}^{\infty}\mu B'_n\le\min(\bover12\epsilon,\beta_r\epsilon^r)$; 
semuanya tentu berjari-jari paling besar $\epsilon$. Misalkan 
$\sequencen{B_n}$ suatu barisan yang menggabungkan $\Cal I_0$ dengan 
$\sequencen{B'_n}$; maka $A\subseteq\bigcup_{n\in\Bbb N}B_n$, setiap $B_n$ 
berjari-jari paling besar $\epsilon$ dan 
      
\Centerline{$\sum_{n=0}^{\infty}\mu B_n 
=\sum_{B\in\Cal I_0}\mu B+\sum_{n=0}^{\infty}\mu B'_n 
\le\mu G+\Bover12\epsilon 
\le\mu^*A+\epsilon$,} 
      
\noindent sedangkan $B_n$ yang pusatnya tidak terletak dalam $A$ harus 
berasal dari barisan $\sequencen{B'_n}$, sehingga jumlah ukurannya paling 
besar $\bover12\epsilon\le\epsilon$. 
}%end of proof of 261F 
      
\cmmnt{\medskip 
      
\noindent{\bf Catatan} Sebenarnya kita dapat (jika $A$ tidak kosong) mengatur 
agar pusat {\it setiap} $B_n$ termasuk dalam $A$. Ini merupakan akibat mudah 
dari Lema Peliputan Besicovitch (lihat \S472 dalam Jilid 4). 
}%end of comment 
      
\exercises{ 
\leader{261X}{Latihan dasar (a)} Tunjukkan bahwa 261C berlaku untuk 
sebarang fungsi bernilai real $f$ yang terintegralkan secara lokal; khususnya, 
untuk sebarang $f\in\eusm L^p(\mu_D)$ dan sebarang $p\ge 1$, dengan $\mu_D$ 
menyatakan ukuran subruang pada $D$. 
%261C 
      
\spheader 261Xb Tunjukkan bahwa 261C, 261Dc, 261Dd dan 261E berlaku untuk 
fungsi bernilai kompleks $f$. 
%261E 
      
\sqheader 261Xc Ambil tiga cakram pada bidang, masing-masing bersinggungan 
dengan dua cakram lainnya, sehingga ketiganya mengurung suatu daerah terbuka 
$R$ dengan tiga titik lancip. Dalam $R$, misalkan $D$ suatu cakram yang 
menyinggung masing-masing dari ketiga cakram semula, dan misalkan 
$R_0$, $R_1$, $R_2$ ketiga komponen $R\setminus D$. Dalam setiap $R_j$, 
misalkan $D_j$ suatu cakram yang menyinggung setiap cakram yang membatasi 
$R_j$, dan misalkan $R_{j0}$, $R_{j1}$, $R_{j2}$ ketiga komponen 
$R_j\setminus D_j$. Lanjutkan proses ini sehingga diperoleh 27 daerah pada 
langkah berikutnya, 81 daerah pada langkah sesudahnya, dan seterusnya. 
      
Tunjukkan bahwa luas total daerah-daerah sisa konvergen ke nol ketika proses 
ini diteruskan tanpa batas. \Hint{bandingkan dengan proses dalam pembuktian 
261B.} 
%261B 
      
\leader{261Y}{Latihan lanjutan (a)} 
%\spheader 261Ya 
Rumuskan suatu definisi abstrak untuk `peliput Vitali', yang berarti suatu 
keluarga himpunan yang memenuhi kesimpulan 261B dalam suatu pengertian, serta 
generalisasi 261C-261E yang bersesuaian dan mencakup (setidaknya) (b)-(d) di 
bawah ini. 
      
\header{261Yb}{\bf (b)} Untuk $x\in\BbbR^r$, $k\in\Bbb N$, misalkan 
$C(x,k)$ adalah kubus setengah terbuka berbentuk 
$\coint{2^{-k}z,2^{-k}(z+\tbf{1})}$, dengan $z\in \Bbb Z^r$, yang memuat 
$x$. Tunjukkan bahwa jika $f$ suatu fungsi yang terintegralkan pada 
$\BbbR^r$, maka 
      
\Centerline{$\lim_{k\to\infty}2^{kr}\int_{C(x,k)}f=f(x)$} 
      
\noindent untuk hampir setiap $x\in\BbbR^r$. 
      
\spheader 261Yc Misalkan $f$ suatu fungsi bernilai real yang terintegralkan 
pada $\BbbR^r$. Tunjukkan bahwa 
      
\Centerline{$\lim_{\delta\downarrow 
0}\Bover{1}{\delta^r}\int_{\coint{x,x+\delta\tbf{1}}}f=f(x)$} 
      
\noindent untuk hampir setiap $x\in\BbbR^r$. 
      
\spheader 261Yd Bekali $X=\{0,1\}^{\Bbb N}$ dengan ukuran lazim $\nu$ 
(254J). Untuk $x\in X$, $k\in\Bbb N$, tetapkan 
$C(x,k)=\{y:y\in X,\,y(i)=x(i)$ untuk $i<k\}$. Tunjukkan bahwa jika $f$ 
sebarang fungsi bernilai real yang terintegralkan pada $X$, maka 
$\lim_{k\to\infty}2^k\int_{C(x,k)}fd\nu=f(x)$ dan 
$\lim_{k\to\infty}2^k\int_{C(x,k)}|f(y)-f(x)|\nu(dy)=0$ untuk hampir setiap 
$x\in X$. 
      
\spheader 261Ye Misalkan $f$ suatu fungsi bernilai real yang terintegralkan 
pada $\BbbR^r$, dan $x$ suatu titik dalam himpunan Lebesgue dari $f$. 
Tunjukkan bahwa untuk setiap $\epsilon>0$ terdapat $\delta>0$ sedemikian 
sehingga 
\discrcenter{390pt}{$|f(x)-\int f(x-y)g(\|y\|)dy|\le\epsilon$ }setiap kali 
$g:\coint{0,\infty}\to\coint{0,\infty}$ suatu fungsi tak menaik sedemikian 
sehingga $\int_{\Bbb R^r}g(\|y\|)dy=1$ dan 
$\int_{B(\tbf{0},\delta)}g(\|y\|)dy\ge 1-\delta$. \Hint{223Yg.} 
      
\spheader 261Yf Misalkan $\frak T$ keluarga himpunan terukur 
$G\subseteq\BbbR^r$ sedemikian sehingga 
\discrcenter{390pt}{$\lim_{\delta\downarrow 0} 
\Bover{\mu(G\cap B(x,\delta))}{\mu B(x,\delta)}=1$ }untuk setiap $x\in G$. 
Tunjukkan bahwa $\frak T$ merupakan suatu topologi pada $\BbbR^r$, yaitu 
{\bf topologi kerapatan} dari $\BbbR^r$. Tunjukkan bahwa suatu fungsi 
$f:\BbbR^r\to\Bbb R$ terukur jika dan hanya jika fungsi itu kontinu dalam 
topologi $\frak T$ pada hampir setiap titik dalam $\Bbb R^r$. 
      
\spheader 261Yg Suatu himpunan $A\subseteq\BbbR^r$ dikatakan {\bf berpori} 
di $x\in\BbbR^r$ jika 
$\limsup_{y\to x}\Bover{\rho(y,A)}{\|y-x\|}>0$, dengan 
$\rho(y,A)=\inf_{z\in A}\|y-z\|$ (atau $\infty$ jika $A$ kosong). 
(i) Tunjukkan bahwa jika $A$ berpori pada semua titiknya, maka himpunan itu 
terabaikan. 
(ii) Tunjukkan bahwa dalam konstruksi 261B, himpunan sisa 
$A\setminus\bigcup\Cal I_0$ berpori pada semua titiknya. 
 
\spheader 261Yh Misalkan $A\subseteq\BbbR^r$ suatu himpunan terbatas dan 
$\Cal I$ suatu keluarga tak kosong dari bola tertutup nontrivial yang meliput 
$A$. Tunjukkan bahwa untuk setiap $\epsilon>0$ terdapat bola-bola yang saling 
lepas $B_0,\ldots,B_n\in\Cal I$ sedemikian sehingga 
$\mu^*A\le(3+\epsilon)^r\sum_{k=0}^n\mu B_k$. 
      
\spheader 261Yi Misalkan $(X,\rho)$ suatu ruang metrik, 
$A\subseteq X$ sebarang himpunan, dan 
$x\mapsto\delta_x:A\to\coint{0,\infty}$ sebarang fungsi terbatas. 
Tunjukkan bahwa jika $\gamma>3$, maka terdapat $A'\subseteq A$ sedemikian 
sehingga (i) $\rho(x,y)>\delta_x+\delta_y$ untuk semua $x$, $y\in A'$ yang 
berbeda; (ii) 
$\bigcup_{x\in A}B(x,\delta_x) 
\subseteq\bigcup_{x\in A'}B(x,\gamma\delta_x)$, dengan $B(x,\alpha)$ 
menyatakan bola tertutup $\{y:\rho(y,x)\le\alpha\}$. 
      
\spheader 261Yj(i) Misalkan $\Cal C$ keluarga himpunan terukur 
$C\subseteq\BbbR^r$ sedemikian sehingga 
\discrcenter{390pt}{$\limsup_{\delta\downarrow 0} 
\Bover{\mu(C\cap B(x,\delta))}{\mu B(x,\delta)}>0$ }untuk setiap $x\in C$. 
Tunjukkan bahwa $\bigcup\Cal C_0\in\Cal C$ untuk setiap 
$\Cal C_0\subseteq\Cal C$. \Hint{215B(iv).}   
(ii) Tunjukkan bahwa setiap gabungan bola tertutup nontrivial dalam 
$\BbbR^r$ terukur Lebesgue. (Bandingkan 415Ye dalam Jilid 4.) 
      
\spheader 261Yk Andaikan bahwa $A\subseteq\BbbR^r$ dan bahwa $\Cal I$ 
merupakan suatu keluarga subhimpunan tertutup dari $\BbbR^r$ sedemikian 
sehingga 
      
\inset{\noindent untuk setiap $x\in A$ terdapat $\eta>0$ sedemikian sehingga 
untuk setiap $\epsilon>0$ terdapat $I\in\Cal I$ sedemikian sehingga $x\in I$ 
dan $0<\eta(\diam I)^r\le\mu I\le\epsilon$.} 
      
\noindent Tunjukkan bahwa terdapat suatu himpunan terhitung yang saling lepas 
$\Cal I_0\subseteq\Cal I$ sedemikian sehingga 
$A\setminus\bigcup\Cal I_0$ terabaikan. 
%261B 
 
\spheader 261Yl Misalkan $\frak T'$ keluarga himpunan terukur 
$G\subseteq\BbbR^r$ sedemikian sehingga setiap kali $x\in G$ dan 
$\epsilon>0$, terdapat $\delta>0$ sedemikian sehingga 
$\mu(G\cap I)\ge(1-\epsilon)\mu I$ setiap kali $I$ suatu interval yang memuat 
$x$ dan termuat dalam $B(x,\delta)$. Tunjukkan bahwa 
$\frak T'$ merupakan suatu topologi pada $\BbbR^r$ yang berada di antara 
topologi kerapatan (261Yf) dan topologi Euklides. 
}%end of exercises 
      
\endnotes{ 
\Notesheader{261} 
Dalam pembuktian 261B-261E di atas, saya telah berusaha sebaik mungkin untuk 
mengikuti alur kasus satu dimensi; bagian ini merupakan serangkaian 
generalisasi dari pekerjaan dalam \S\S221 dan 223. 
      
Jelas bahwa gagasan 261A/261B dapat digunakan pada bentuk selain bola. Agar 
gagasan itu bekerja dalam bentuk di atas, kita memerlukan suatu keluarga 
$\Cal I$ sedemikian sehingga terdapat konstanta $K$ yang memenuhi 
      
\Centerline{$\mu I'\le K\mu I$} 
      
\noindent untuk setiap $I\in \Cal I$, dengan kita menuliskan 
      
\Centerline{$I'=\{x:\inf_{y\in I}\|x-y\|\le\diam(I)\}$.} 
      
\noindent Jelas bahwa hal ini berlaku bagi banyak kelas $\Cal I$ yang 
ditentukan oleh bentuk himpunan-himpunan yang terlibat; misalnya, jika 
$E\subseteq\BbbR^r$ sebarang himpunan terbatas dengan ukuran positif, 
keluarga $\Cal I=\{x+\delta E:x\in\BbbR^r,\,\delta>0\}$ memenuhi syarat itu. 
      
Dalam 261Ya saya menantang Anda untuk menemukan generalisasi yang sesuai bagi 
argumen-argumen yang bergantung pada kesimpulan 261B. 
      
Cara lain menggunakan 261B adalah mengatakan bahwa karena himpunan secara 
esensial dapat diliput oleh barisan bola yang {\it saling lepas}, seharusnya 
bola, bukan interval setengah terbuka, dapat digunakan dalam definisi ukuran 
Lebesgue pada $\BbbR^r$. Hal ini memang benar (261F). Kesulitan menggunakan 
bola dalam definisi dasar muncul tepat pada permulaan, ketika membuktikan 
bahwa jika suatu bola diliput oleh berhingga banyak bola, maka jumlah volume 
bola-bola peliput itu setidaknya sebesar volume bola yang diliput. (Ada suatu 
siasat, dengan menggunakan kekompakan bola tertutup dan keterbukaan bola 
terbuka, untuk memperluas pembuktian demikian ke peliput tak hingga.) Tentu 
saja Anda dapat menganggap fakta ini `elementer', atas dasar bahwa Archimedes 
tentu akan menyadari seandainya fakta itu tidak benar, tetapi bagaimanapun 
fakta itu akan cukup menantang untuk dibuktikan, kecuali Anda bersedia 
menunggu suatu versi teorema Fubini, sebagaimana dilakukan beberapa penulis. 
      
Saya memberikan hasil-hasil dalam 261C-261D untuk subhimpunan sebarang $D$ 
dari $\Bbb R^r$, bukan karena saya membayangkan penerapan yang menjadikan 
subhimpunan tak terukur penting, melainkan karena saya ingin memungkinkan 
kita memperhatikan kapan keterukuran berperan. Tentu saja integral 
$\int_Dfd\mu$ perlu ditafsirkan dengan cara yang ditetapkan dalam \S214. 
Rahasia argumennya terungkap dalam bagian (c) pembuktian 261C, tempat saya 
mengandalkan fakta bahwa jika $f$ terintegralkan pada $D$, maka terdapat 
suatu fungsi terintegralkan $\tilde f:\BbbR^r\to\Bbb R$ sedemikian sehingga 
$\int_F\tilde f=\int_{D\cap F}f$ untuk setiap himpunan terukur 
$F\subseteq\BbbR^r$. Pada hakikatnya, untuk semua persoalan yang dibahas di 
sini, kita dapat mengganti $f$, $D$ dengan $\tilde f$, $\BbbR^r$. 
      
Gagasan 261C ialah bahwa, untuk hampir setiap $x$, $f(x)$ didekati oleh nilai 
rata-ratanya pada bola kecil $B(x,\delta)$, dengan mengabaikan nilai-nilai 
yang hilang pada $B(x,\delta)\setminus(D\cap\dom f)$; 261E merupakan versi 
lebih tajam dari gagasan yang sama. Rumus-rumus 261C-261E terutama memuat 
ungkapan $\mu B(x,\delta)$. Tentu saja ini hanyalah $\beta_r\delta^r$. Namun, 
menurut saya membiarkannya tanpa dikembangkan justru lebih memperjelas, 
sekaligus menghindari indeks bawah dan atas, karena dengan demikian lebih 
jelas apa pokok sesungguhnya dari teorema-teorema kerapatan ini. Dalam \S472 
dari Jilid 4 saya akan kembali ke materi ini dan menunjukkan bahwa sebagian 
besar gagasannya, dalam kadar yang mengejutkan, dapat diterapkan pada 
sebarang ukuran Radon pada $\BbbR^r$, walaupun teorema Vitali (dalam bentuk 
yang dinyatakan di sini) tidak lagi berlaku. 
      
}%end of notes 
      
      
\discrpage 
      
